\documentclass{../yalumba}

\usepackage{multirow}
\usepackage{array}

\title{The Representation Bottleneck:\\A Perturbation Audit of Module-Graph\\Features for Genomic Ecosystem Comparison}
\author{Ajax Davis}
\date{April 2026}

\setabstract{%
We present the results of a systematic perturbation audit of the
module-graph representation that has served as the computational substrate
for ten versions of Spectral Ecology and three prior symbiogenesis
algorithms.  This is a \textbf{negative result paper}.  The audit applies
three complementary stress tests---a coverage ladder (subsampling at
100\%, 80\%, 60\%, and 40\% of 50,000 reads), cross-sample variance
analysis, and a read-order shuffle test---to the 15 features extracted
from k-mer co-occurrence graphs constructed for six members of the CEPH
1463 pedigree.  The findings are unambiguous and sobering.  First, only
5 of 15 features achieve self-perturbation stability (coefficient of
variation below 0.15 across coverage levels): module count, edge density,
mean degree, mean curvature, and median curvature.  The remaining 10
features, including all five ecological role counts that several
algorithm versions weighted heavily, are \emph{unstable under
subsampling within the same sample}.  Second, the signal-to-nuisance
ratio---defined as cross-sample variance divided by self-perturbation
variance---exceeds 3 for \textbf{zero} of 15 features.  The best ratio
is 2.93 (edge density); the worst is 0.59 (scaffold role count).  Four
features, including two of the three curvature statistics that Spectral
Ecology v9 scored, have ratios below 1.0, meaning that variation
introduced by subsampling the \emph{same} genome exceeds variation
between \emph{different} genomes.  For these features, the
representation is literally noisier than the signal it is supposed to
carry.  Third, and most devastatingly, shuffling read order---a
transformation that changes zero biological content---produces
catastrophic changes in graph structure: module count increases by 45\%
to 184\%, triangle count increases by 2\% to 607\%, and edge count
increases by 27\% to 264\%.  The module-graph representation is not
merely noisy; it is \emph{non-deterministic} with respect to read
ordering, a property that no downstream invariant can correct.  These
results explain why the ten-version Spectral Ecology arc, despite
producing genuine theoretical advances in coherence fields, intrinsic
curvature distributions, holonomy deficits, and Sinkhorn-regularized
optimal transport, never achieved a positive weakest-pair gap: the
scoring algorithms were optimizing invariants on an unstable
representation.  The downstream sophistication was outrunning
upstream identifiability.  We provide detailed tables of all 15 features
across all three tests, analyze the two ``least-bad'' features (edge
density at ratio 2.93 and median curvature at ratio 2.47) as potential
starting points for a stabilized representation, and offer concrete
recommendations for the next phase: deterministic module extraction,
coverage-matched subsampling, rank-based normalization, and an acceptance
criterion requiring signal-to-nuisance ratio above 3 for every scored
feature before any new algorithm version is attempted.  This paper
changes the direction of the yalumba research program.  It is, we argue,
the most important experiment the project has conducted.}

\begin{document}
\maketitle

\tableofcontents
\newpage


% ════════════════════════════════════════════════════════════════════
\section{Introduction}
\label{sec:intro}

\subsection{Why Negative Results Matter}

The history of computational science is littered with methods that
optimized the wrong objective on the wrong substrate.  In machine
learning, feature engineering on unstable representations produced years
of incremental papers before practitioners learned to audit the
representation itself~\cite{bendavid2006stability}.  In genomics, early
reference-free methods achieved impressive-looking metrics on benchmarks
whose variance properties were never characterized~\cite{fan2021reference}.
In ecology, decades of community comparison used indices whose sensitivity
to sampling effort exceeded their sensitivity to actual community
differences~\cite{bray1957, morisita1959}.

\begin{pullquote}
The most important experiment in a research program is not the one that
succeeds.  It is the one that reveals what success was being measured
against.
\end{pullquote}

This paper reports such an experiment.  After ten versions of Spectral
Ecology---a research arc spanning sparse graphs, dense co-occurrence
graphs, optimal transport, coherence fields, persistence fields,
Sinkhorn regularization, holonomy deficits, intrinsic curvature, and
hybrid invariants---we step back and ask a question that should have been
asked before version~1:

\begin{keyfinding}
Is the module-graph representation stable enough for \emph{any}
downstream invariant to distinguish biological signal from
representation noise?
\end{keyfinding}

The answer is no.

\subsection{The Ten-Version Arc as Context}

Table~\ref{tab:arc-summary} summarizes the Spectral Ecology program.
Each version introduced increasingly sophisticated mathematical
machinery: Laplacian eigendecomposition (v1--v2), Wasserstein transport
(v3), coherence fields (v4--v5), persistence fields (v6), Sinkhorn
regularization with cooperative weighting (v7), gauge-theoretic holonomy
(v8), Ollivier--Ricci curvature distributions (v9), and hybrid
intrinsic--relational composites (v10).

\begin{table}[htbp]
\centering
\tablestyle
\caption{The Spectral Ecology arc: ten versions, ten paradigms.}
\label{tab:arc-summary}
\begin{tabular}{@{}llrr@{}}
\toprule
\rowcolor{table-header}
\textcolor{white}{\textbf{Version}} &
\textcolor{white}{\textbf{Core Idea}} &
\textcolor{white}{\textbf{Sep.}} &
\textcolor{white}{\textbf{Gap}} \\
\midrule
v1 & Sparse adjacency, eigenvalue density  & $-0.74\%$ & $-26.87\%$ \\
v2 & Dense co-occurrence, heat kernel      & $+1.55\%$ & $-22.53\%$ \\
v3 & Wasserstein transport                 & $+4.93\%$ & $-54.67\%$ \\
v4 & Coherence fields                      & $+2.10\%$ & $-18.40\%$ \\
v5 & Multi-scale coherence                 & $+1.80\%$ & $-12.30\%$ \\
v6 & Persistence fields                    & $+2.40\%$ & $-3.21\%$ \\
v7 & Sinkhorn cooperative stability        & $+0.87\%$ & $-0.61\%$ \\
v8 & Holonomy gauge invariants             & $+1.62\%$ & $-6.41\%$ \\
v9 & Intrinsic curvature distributions     & $+3.20\%$ & $-7.93\%$ \\
v10 & Hybrid intrinsic--relational         & $+1.50\%$ & $-1.66\%$ \\
\bottomrule
\end{tabular}
\end{table}

\ymarginnote{The gap trajectory: $-26.87\% \to -0.61\% \to -1.66\%$.
Progress is real but the gap never turns positive.}

The trajectory is suggestive.  Separation improved from $-0.74\%$ to
$+4.93\%$ (v3), and the weakest gap narrowed from $-26.87\%$ to
$-0.61\%$ (v7).  But the gap never turned positive.  The persistent
failure mode was the spouse pair (Father--Mother), which consistently
scored among the top-ranked pairs despite being unrelated, due to
similar sequencing coverage producing similar module graph structure.

Every version attributed this to a \emph{scoring confound}---coverage
similarity inflating pairwise scores---and attempted to fix it
downstream: null-model deconfounding (v2), self-baseline normalization
(v8), distributional comparison (v9).  None of these attempts succeeded
in closing the gap to positive.

A pattern emerges from the arc.  The three eras of the program---within-sample
spectral summaries (v1--v2), cross-sample relational methods (v3--v6), and
intrinsic geometric invariants (v7--v10)---each attacked the \emph{scoring}
problem with increasing mathematical sophistication.  But none attacked the
\emph{representation} problem.  The module graph was accepted as given, like a
microscope that nobody thought to calibrate.

This paper provides the explanation: \textbf{the confound is not in the
scoring.  It is in the representation.}

\subsection{The Feedback That Prompted This Audit}

The impetus for this experiment was a piece of external feedback on the
theoretical framework paper~\cite{nakahara2003geometry}: ``downstream
sophistication is outrunning upstream identifiability.''  The observation
was that the program had invested enormous effort in invariants
(eigenvalue densities, heat kernel traces, curvature distributions,
transport plans, holonomy deficits) while making essentially zero
investment in understanding whether the input to those invariants---the
module graph itself---was a reliable measurement instrument.

\begin{definition}
A \textbf{measurement instrument} is reliable if repeated measurements
of the same object yield similar results.  In our context, the ``object''
is a genome, the ``measurement'' is the module-graph representation, and
``repeated measurements'' are obtained by subsampling reads or
reordering them.  A reliable representation has low self-perturbation
variance relative to cross-sample variance.
\end{definition}

This is not a novel concept.  It is the foundation of measurement
science: an instrument must have lower measurement error than the effect
size it is trying to detect.  But in ten versions of algorithm
development, this property was never tested.

\subsection{Why This Is the Most Important Experiment}

The module-graph representation underlies every algorithm in the
symbiogenesis program.  Module Persistence~v3 counts shared modules.
Coalition Transfer~v3 counts shared coalitions.  All ten versions of
Spectral Ecology extract invariants from module co-occurrence graphs.
If the representation is unreliable, then:

\begin{itemize}
  \item Algorithms that pass (Module Persistence, Coalition Transfer)
    may be exploiting the \emph{stable subset} of features by accident.
  \item Algorithms that fail (Module Stability, Spectral Ecology) may
    be scoring features whose variation is dominated by representation
    noise.
  \item The persistent coverage confound may not be a confound at all,
    but a consequence of the representation collapsing biological
    diversity into coverage-dependent structure.
\end{itemize}

This experiment tests all three hypotheses simultaneously.

\subsection{Relationship to Passing Algorithms}

Module Persistence v3 ($+3.36\%$ separation, $+0.13\%$ gap) and
Coalition Transfer v3 ($+1.97\%$ separation, $+0.06\%$ gap) both pass
the kinship detection threshold.  If the representation is unreliable,
why do they work?

The answer, previewed here and explored in Section~\ref{sec:discussion},
is that these algorithms happen to exploit the stable subset of
features.  Module Persistence scores the overlap of \emph{which modules
exist}---a property closely related to module count, which has a CV of
0.11.  Coalition Transfer scores co-occurrence of module pairs---a
property related to edge density, which has a CV of 0.07.  Neither
algorithm directly scores triangle counts, role distributions, or
curvature statistics.  Their success is not evidence that the
representation is sound.  It is evidence that a few features within the
representation carry enough signal for simple algorithms---and that
sophisticated algorithms that score the full feature set are worse off,
not better off, than simple ones that accidentally restrict themselves
to the stable features.

This is the cruelest finding of the audit: \textbf{the simplest
algorithms in the project were the most robust, and the most
sophisticated were the most fragile.}


% ════════════════════════════════════════════════════════════════════
\section{Background}
\label{sec:background}

\subsection{The Module-Graph Representation}
\label{sec:bg-representation}

The module-graph representation is constructed as follows.  Given a set
of reads from a FASTQ file:

\begin{enumerate}
  \item \textbf{K-mer extraction.}  Canonical k-mers (lexicographically
    smaller of k-mer and its reverse complement) are extracted from each
    read using a rolling hash~\cite{karp1987rabin}.  Typical $k = 21$.
  \item \textbf{Frequency filtering.}  K-mers are filtered by frequency
    to remove sequencing errors (low count) and repetitive elements
    (high count).  Typical thresholds: frequency $\geq 2$, frequency
    $\leq 100$.
  \item \textbf{Co-occurrence counting.}  For each pair of k-mers
    appearing on the same read within a window, a co-occurrence count is
    incremented.  This produces a symmetric co-occurrence matrix.
  \item \textbf{Graph construction.}  K-mers become nodes.  Edges are
    created between k-mers whose co-occurrence exceeds a threshold
    (typically the mean co-occurrence plus one standard deviation).
  \item \textbf{Module extraction.}  Connected components of the graph
    are identified as ``modules.''  Each module is a set of k-mers that
    co-occur frequently on reads.
  \item \textbf{Feature extraction.}  15 features are computed from the
    module graph: module count, triangle count, edge count, edge density,
    mean/max degree, curvature statistics (mean, std, p50, p90), and
    ecological role counts (core, satellite, bridge, scaffold, boundary).
\end{enumerate}

\subsection{What Ten Versions Assumed}

Every version of Spectral Ecology assumed, implicitly or explicitly,
that the module graph captured something biologically meaningful and
reproducible about the genome.  Specifically:

\begin{itemize}
  \item \textbf{v1--v2} assumed that eigenvalue densities of the graph
    Laplacian were genome-specific signatures.
  \item \textbf{v3} assumed that the graph structure was stable enough
    for optimal transport between graphs to be meaningful.
  \item \textbf{v4--v6} assumed that multi-scale heat kernel signatures
    encoded persistent structural features.
  \item \textbf{v7} assumed that cooperative stability (agreement across
    multiple invariants) would cancel noise.
  \item \textbf{v8} assumed that holonomy---path-dependent parallel
    transport on the graph---was a gauge-invariant property of the
    biology, not an artifact of the representation.
  \item \textbf{v9} assumed that Ollivier--Ricci curvature distributions
    captured intrinsic geometric properties that were independent of
    graph size and density.
  \item \textbf{v10} assumed that combining intrinsic and relational
    invariants would suppress each other's failure modes.
\end{itemize}

None of these assumptions were tested against the null hypothesis that
the module graph's structure is dominated by coverage-dependent and
order-dependent artifacts.

\ymarginnote{Ten versions of assumptions.  Zero versions of validation.
This audit is the first.}

The implicit chain of reasoning was: ``the graph captures biology
$\Rightarrow$ invariants of the graph capture invariants of the biology
$\Rightarrow$ comparing invariants compares biology.''  The first link
in this chain---the graph captures biology---was never validated.  If
it is false, the rest of the chain is vacuous regardless of the
mathematical sophistication applied downstream.

\subsection{The Concept of Signal-to-Nuisance Ratio}

We borrow from signal processing the idea that a measurement system's
utility depends not on the absolute magnitude of signal it detects, but
on the ratio of signal to noise.  In our context:

\begin{definition}
The \textbf{signal-to-nuisance ratio} (SNR) of a feature is:
\begin{equation}
  \text{SNR} = \frac{\text{CV}_{\text{cross-sample}}}
                     {\text{CV}_{\text{self-perturbation}}}
  \label{eq:snr}
\end{equation}
where $\text{CV}$ is the coefficient of variation.  Cross-sample CV
measures how much the feature varies between different genomes at full
coverage.  Self-perturbation CV measures how much the feature varies
when the \emph{same} genome is measured at different coverage levels.
\end{definition}

An SNR of 1.0 means the feature varies as much between measurements
of the same genome as it does between different genomes.  Such a feature
carries no usable signal.  An SNR of 3.0 means cross-sample variation
is three times self-perturbation variation---a minimal threshold for
the feature to contribute to discrimination.

\begin{limitation}
This definition of SNR uses the coefficient of variation, which assumes
roughly linear scaling of standard deviation with mean.  For features
with heavy-tailed distributions or zero-inflated counts (such as bridge
and scaffold role counts), the CV may overstate instability.  We use
it here as a simple, interpretable metric; more sophisticated measures
(intraclass correlation, measurement invariance models) are left for
future work.
\end{limitation}

\subsection{Why Perturbation Testing Matters}

Perturbation testing is not algorithm design.  It is \emph{measurement
science}~\cite{efron1979bootstrap, hennig2007cluster}.  The question is
not ``can we build a clever algorithm?'' but ``is the input to that
algorithm a reliable instrument?''

\ymarginnote{Metrology---the science of measurement---requires that
instrument precision be characterized before any measurement is
reported.  We never did this.}

In classical genomics, perturbation testing is routine.  Variant callers
are evaluated by subsampling reads and measuring concordance.  IBD
detectors are tested across coverage
levels~\cite{browning2012ibd, browning2013refinedibd}.
Reference-free methods validate stability under read
subsets~\cite{ondov2016mash, brinda2023phylogenetic}.

The yalumba project skipped this step.  For ten versions, the
representation was treated as a given---a fixed input whose properties
were unquestioned.  The audit reported here is, in retrospect, the
experiment that should have been conducted before version~1.

The analogy to ecology is instructive.  When Bray and Curtis introduced
their dissimilarity index in 1957~\cite{bray1957}, they accompanied it
with extensive discussion of sampling effort and its effect on community
representation.  Morisita's overlap index~\cite{morisita1959} was
explicitly designed to be robust to sample size.  The ecological
literature understood, sixty years ago, that \emph{you cannot compare
communities if your representation of each community depends on how hard
you looked}.  The genomics field sometimes forgets this lesson.


% ════════════════════════════════════════════════════════════════════
\section{Methods}
\label{sec:methods}

\subsection{Dataset}

All experiments use the CEPH 1463 pedigree~\cite{ceph1463}, a
three-generation family:

\begin{pedigree}
\begin{verbatim}
       NA12891 (Pat. GF) --- NA12892 (Pat. GM)
                    |
                NA12877 (Father)
                    |
                    +---- NA12878 (Mother)
                    |
       NA12889 (Mat. GF) --- NA12890 (Mat. GM)
\end{verbatim}
\end{pedigree}

Six samples, each processed from Illumina short-read FASTQ files.
50,000 reads per sample at full coverage.  Related pairs: Father--Mother
is the \emph{spouse pair} (unrelated).  Father--Pat.~GF, Father--Pat.~GM,
Mother--Mat.~GF, Mother--Mat.~GM are the four parent--child pairs
(related).

\subsection{Coverage Ladder}
\label{sec:coverage-ladder}

For each of the six samples, we construct module graphs at four coverage
levels by subsampling reads:

\begin{itemize}
  \item \textbf{100\%:} all 50,000 reads
  \item \textbf{80\%:} 40,000 reads (random subsample, fixed seed)
  \item \textbf{60\%:} 30,000 reads (random subsample, fixed seed)
  \item \textbf{40\%:} 20,000 reads (random subsample, fixed seed)
\end{itemize}

Each coverage level produces an independent module graph.  Features are
extracted from each graph.  The coefficient of variation across the four
coverage levels, for each feature within each sample, measures
self-perturbation instability.

\subsection{Shuffle Test}
\label{sec:shuffle-test}

For each of the six samples, we take the full 50,000 reads and
\textbf{randomize their order} using a deterministic permutation (fixed
seed).  The shuffled reads contain \emph{exactly the same biological
content}---the same k-mers, the same k-mer frequencies, the same
sequences.  The only difference is the order in which reads are
processed.

If the module-graph representation were a function of the biological
content alone, the shuffle test would produce identical graphs.  Any
change in graph structure is pure representation artifact.

\subsection{Feature Extraction}

From each module graph, we extract 15 features organized into four
categories:

\begin{enumerate}
  \item \textbf{Graph topology (4 features):} module count, triangle
    count, edge count, edge density.
  \item \textbf{Degree statistics (2 features):} mean degree, max degree.
  \item \textbf{Curvature statistics (4 features):} Ollivier--Ricci
    curvature mean, standard deviation, 50th percentile (p50), and 90th
    percentile (p90).
  \item \textbf{Ecological role counts (5 features):} core, satellite,
    bridge, scaffold, boundary.  These are assigned by thresholding degree,
    betweenness centrality, and local clustering
    coefficient~\cite{hartwell1999cell, milo2002network}.
\end{enumerate}

\subsection{Self-Perturbation CV}

For each feature $f$ and each sample $s$, we compute the coefficient
of variation across the four coverage levels:

\begin{equation}
  \text{CV}_{\text{self}}(f, s) =
    \frac{\sigma_{f,s}}{\mu_{f,s}}
  \label{eq:cv-self}
\end{equation}

where $\mu_{f,s}$ and $\sigma_{f,s}$ are the mean and standard
deviation of feature $f$ across the four coverage levels for sample $s$.
The reported self-perturbation CV is the mean across all six samples:

\begin{equation}
  \overline{\text{CV}}_{\text{self}}(f) =
    \frac{1}{6} \sum_{s=1}^{6} \text{CV}_{\text{self}}(f, s)
  \label{eq:cv-self-mean}
\end{equation}

\subsection{Cross-Sample CV}

At 100\% coverage, we compute the coefficient of variation of each
feature across the six samples:

\begin{equation}
  \text{CV}_{\text{cross}}(f) =
    \frac{\sigma_f}{\mu_f}
  \label{eq:cv-cross}
\end{equation}

This measures the variation that exists between genuinely different
genomes under identical processing conditions.

\subsection{Signal-to-Nuisance Ratio}

The ratio is computed per Equation~\ref{eq:snr}.  Features with
$\text{SNR} < 1$ are dominated by self-perturbation noise.  Features
with $1 \leq \text{SNR} < 3$ are marginal.  Features with
$\text{SNR} \geq 3$ are minimally acceptable for downstream scoring.

\begin{keyfinding}
The acceptance criterion is $\text{SNR} \geq 3$.  This threshold is
deliberately conservative.  In psychometrics and measurement theory,
an instrument with reliability below 0.7 (roughly corresponding to
$\text{SNR} < 3$) is considered unfit for individual-level
inference~\cite{hennig2007cluster}.  We adopt the same standard.
\end{keyfinding}

\subsection{Experimental Protocol}

All experiments use identical k-mer parameters ($k = 21$, frequency
filter $\geq 2$, $\leq 100$), identical graph construction thresholds,
and identical feature extraction code.  The only variables are:

\begin{enumerate}
  \item \textbf{Coverage ladder:} number of reads used (100\%, 80\%,
    60\%, 40\%).
  \item \textbf{Shuffle test:} order of reads (original vs. permuted).
\end{enumerate}

All random operations use fixed seeds for reproducibility.

\subsection{Stability Classification}

We classify features into three stability tiers based on self-perturbation
CV:

\begin{definition}
\begin{itemize}
  \item \textbf{STABLE:} Mean self-perturbation CV $< 0.15$.  The
    feature changes by less than 15\% on average when coverage is reduced
    by up to 60\%.
  \item \textbf{MODERATE:} Mean self-perturbation CV in $[0.15, 0.40)$.
    The feature changes by 15--40\% under coverage perturbation.
  \item \textbf{UNSTABLE:} Mean self-perturbation CV $\geq 0.40$.  The
    feature changes by more than 40\%, or its standard deviation exceeds
    its mean at some coverage levels.
\end{itemize}
\end{definition}

These thresholds are deliberately generous.  A 15\% CV is already
concerning for a measurement instrument---it means that measuring the
same genome at 80\% coverage instead of 100\% can change the feature
by more than one standard deviation.  We adopt lenient thresholds to
ensure that features classified as UNSTABLE are unambiguously
unreliable.

\subsection{Verdict Classification for SNR}

For the signal-to-nuisance ratio, we use three verdicts:

\begin{itemize}
  \item \textbf{NOISE $>$ SIGNAL} ($\text{SNR} < 1$): Self-perturbation
    variance exceeds cross-sample variance.  The feature cannot
    distinguish genomes.
  \item \textbf{MARGINAL} ($1 \leq \text{SNR} < 3$): Cross-sample
    variance exceeds self-perturbation variance, but not by enough
    for reliable discrimination.
  \item \textbf{ACCEPTABLE} ($\text{SNR} \geq 3$): The feature has
    enough signal relative to noise for downstream scoring.
\end{itemize}

As the results will show, no feature reaches the ACCEPTABLE tier.


% ════════════════════════════════════════════════════════════════════
\section{Results}
\label{sec:results}

\subsection{Self-Perturbation Analysis}
\label{sec:self-perturbation}

Table~\ref{tab:self-perturbation} presents the self-perturbation CV
for all 15 features across the six CEPH 1463 samples.

\begin{table}[htbp]
\centering
\tablestyle
\caption{Self-perturbation analysis.  CV = coefficient of variation
across coverage levels (100\%, 80\%, 60\%, 40\%) averaged over six
samples.  Features with mean CV $< 0.15$ are classified as STABLE;
$0.15$--$0.40$ as MODERATE; $> 0.40$ as UNSTABLE.}
\label{tab:self-perturbation}
\begin{tabular}{@{}lrrrl@{}}
\toprule
\rowcolor{table-header}
\textcolor{white}{\textbf{Feature}} &
\textcolor{white}{\textbf{Mean CV}} &
\textcolor{white}{\textbf{Min CV}} &
\textcolor{white}{\textbf{Max CV}} &
\textcolor{white}{\textbf{Verdict}} \\
\midrule
moduleCount    & 0.1059 & 0.0154 & 0.2119 & STABLE \\
triangleCount  & 0.4386 & 0.2244 & 0.8157 & UNSTABLE \\
edgeCount      & 0.2303 & 0.0758 & 0.4579 & MODERATE \\
edgeDensity    & 0.0725 & 0.0147 & 0.1843 & STABLE \\
meanDegree     & 0.1249 & 0.0446 & 0.2675 & STABLE \\
maxDegree      & 0.2947 & 0.0850 & 0.6563 & MODERATE \\
curvMean       & 0.1373 & 0.0274 & 0.3549 & STABLE \\
curvStd        & 0.2452 & 0.0656 & 0.5500 & MODERATE \\
curvP50        & 0.0660 & 0.0168 & 0.1006 & STABLE \\
curvP90        & 0.1949 & 0.0205 & 0.5494 & MODERATE \\
role\_core      & 0.5458 & 0.0831 & 1.5508 & UNSTABLE \\
role\_satellite & 0.4689 & 0.1884 & 0.7971 & UNSTABLE \\
role\_bridge    & 1.0875 & 0.0000 & 1.7321 & UNSTABLE \\
role\_scaffold  & 0.8034 & 0.0000 & 1.7321 & UNSTABLE \\
role\_boundary  & 0.7251 & 0.3333 & 1.7321 & UNSTABLE \\
\bottomrule
\end{tabular}
\end{table}

\begin{keyfinding}
Only 5 of 15 features are stable under coverage perturbation: module
count (CV = 0.106), edge density (CV = 0.073), mean degree (CV = 0.125),
mean curvature (CV = 0.137), and median curvature (CV = 0.066).  All
five ecological role counts are \textbf{unstable}, with CVs ranging from
0.47 (satellite) to 1.09 (bridge).
\end{keyfinding}

The instability of role counts is particularly damaging.  The bridge
role has a maximum CV of 1.73---a value that indicates the standard
deviation across coverage levels \emph{exceeds the mean}.  For several
samples, the bridge count is zero at some coverage levels and nonzero
at others, producing undefined or extreme CVs.  Similarly, scaffold and
boundary roles fluctuate wildly, with max CVs of 1.73.

The moderate features (edge count, max degree, curvature std, curvature
p90) sit in an uncomfortable middle ground.  Their CVs of 0.19--0.29
mean that a 40\% reduction in coverage can change the feature value by
20--30\%.  This is not catastrophic, but it is far from the stability
needed for reliable comparison.

\subsection{Cross-Sample Variance}
\label{sec:cross-sample}

Table~\ref{tab:cross-sample} presents the cross-sample variance at
100\% coverage.

\begin{table}[htbp]
\centering
\tablestyle
\caption{Cross-sample variance at 100\% coverage (6 samples).  These
numbers represent the ``signal'' side of the signal-to-nuisance ratio:
how much do features vary between genuinely different genomes?}
\label{tab:cross-sample}
\begin{tabular}{@{}lrrrr@{}}
\toprule
\rowcolor{table-header}
\textcolor{white}{\textbf{Feature}} &
\textcolor{white}{\textbf{Mean}} &
\textcolor{white}{\textbf{Std}} &
\textcolor{white}{\textbf{CV}} &
\textcolor{white}{\textbf{Range}} \\
\midrule
moduleCount     & 588.67 &   90.39 & 0.1536 & 457--712 \\
triangleCount   & 3305.00 & 1363.32 & 0.4125 & 1256--4867 \\
edgeCount       & 1675.33 &  394.18 & 0.2353 & 1110--2251 \\
edgeDensity     &    0.01 &    0.00 & 0.2124 & 0.004--0.009 \\
meanDegree      &    5.66 &    0.83 & 0.1465 & 4.5--6.4 \\
maxDegree       &   81.00 &   38.39 & 0.4739 & 37--144 \\
curvMean        &    1.22 &    0.15 & 0.1244 & 1.1--1.4 \\
curvStd         &    1.57 &    0.36 & 0.2275 & 1.0--1.9 \\
curvP50         &    0.79 &    0.13 & 0.1633 & 0.6--1.0 \\
curvP90         &    2.49 &    0.68 & 0.2717 & 1.5--3.4 \\
role\_core       &  411.50 &  218.71 & 0.5315 & 14--629 \\
role\_satellite  &   91.33 &   71.49 & 0.7827 & 19--223 \\
role\_bridge     &   76.67 &  170.99 & 2.2302 & 0--459 \\
role\_scaffold   &    7.33 &    3.50 & 0.4767 & 0--11 \\
role\_boundary   &    1.83 &    1.57 & 0.8576 & 0--5 \\
\bottomrule
\end{tabular}
\end{table}

The cross-sample CVs paint a superficially encouraging picture.  Role
counts have high cross-sample variance---bridge role has a CV of 2.23,
suggesting enormous variation between samples.  But high cross-sample
variance is \emph{only meaningful if self-perturbation variance is low}.
A feature that varies wildly between samples \emph{and} varies wildly
within a sample is just noisy, not discriminative.

\ymarginnote{High cross-sample CV is necessary but not sufficient.  The
signal-to-nuisance ratio asks: is the cross-sample variation genuine,
or is it the same noise that appears within samples?}

Note the role count ranges.  Core ranges from 14 to 629---a 45-fold
range.  Bridge ranges from 0 to 459.  Scaffold from 0 to 11.  These
enormous ranges will produce the appearance of discrimination in any
downstream algorithm, but we must verify that they reflect biology
rather than coverage artifacts.

\subsection{Signal-to-Nuisance Ratio}
\label{sec:snr}

Table~\ref{tab:snr} presents the signal-to-nuisance ratio for all 15
features.  This is the central result of the paper.

\begin{table}[htbp]
\centering
\tablestyle
\caption{Signal-to-nuisance ratio (SNR) for all 15 features.
$\text{SNR} = \text{CV}_{\text{cross}} / \text{CV}_{\text{self}}$.
Features with $\text{SNR} < 1$ carry more noise than signal.  No
feature achieves the acceptance threshold of $\text{SNR} \geq 3$.}
\label{tab:snr}
\begin{tabular}{@{}lrrrl@{}}
\toprule
\rowcolor{table-header}
\textcolor{white}{\textbf{Feature}} &
\textcolor{white}{\textbf{Self CV}} &
\textcolor{white}{\textbf{Cross CV}} &
\textcolor{white}{\textbf{Ratio}} &
\textcolor{white}{\textbf{Verdict}} \\
\midrule
moduleCount     & 0.1059 & 0.1536 & 1.45 & MARGINAL \\
triangleCount   & 0.4386 & 0.4125 & 0.94 & NOISE $>$ SIGNAL \\
edgeCount       & 0.2303 & 0.2353 & 1.02 & MARGINAL \\
edgeDensity     & 0.0725 & 0.2124 & 2.93 & MARGINAL (best) \\
meanDegree      & 0.1249 & 0.1465 & 1.17 & MARGINAL \\
maxDegree       & 0.2947 & 0.4739 & 1.61 & MARGINAL \\
curvMean        & 0.1373 & 0.1244 & 0.91 & NOISE $>$ SIGNAL \\
curvStd         & 0.2452 & 0.2275 & 0.93 & NOISE $>$ SIGNAL \\
curvP50         & 0.0660 & 0.1633 & 2.47 & MARGINAL \\
curvP90         & 0.1949 & 0.2717 & 1.39 & MARGINAL \\
role\_core       & 0.5458 & 0.5315 & 0.97 & NOISE $>$ SIGNAL \\
role\_satellite  & 0.4689 & 0.7827 & 1.67 & MARGINAL \\
role\_bridge     & 1.0875 & 2.2302 & 2.05 & MARGINAL \\
role\_scaffold   & 0.8034 & 0.4767 & 0.59 & NOISE $>$ SIGNAL \\
role\_boundary   & 0.7251 & 0.8576 & 1.18 & MARGINAL \\
\bottomrule
\end{tabular}
\end{table}

\begin{keyfinding}
\textbf{Zero features achieve an SNR of 3 or above.}  The best SNR is
2.93 (edge density).  The second-best is 2.47 (median curvature, curvP50).
After that, the ratios drop below 2.1.  Four features have
$\text{SNR} < 1$, meaning self-perturbation noise exceeds cross-sample
variation: triangle count (0.94), mean curvature (0.91), curvature
standard deviation (0.93), core role count (0.97), and scaffold role
count (0.59).
\end{keyfinding}

Figure~\ref{fig:snr-landscape} categorizes the features visually.

\begin{figure}[htbp]
\centering
\begin{tikzpicture}[
  x=4cm, y=4cm,
  marginal/.style={fill=yalumba-gold!30, circle, inner sep=2pt,
    minimum size=5pt},
  noise/.style={fill=yalumba-wine!50, circle, inner sep=2pt,
    minimum size=5pt},
]
% Axes
\draw[->] (0,0) -- (2.6,0) node[right] {\small\sffamily Cross-sample CV};
\draw[->] (0,0) -- (0,1.4) node[above] {\small\sffamily Self-perturbation CV};

% Reference lines
\draw[dashed, yalumba-wine, thick] (0,0) -- (2.4,2.4*0.333)
  node[right, font=\tiny\sffamily] {SNR = 3};
\draw[dashed, yalumba-slate, thin] (0,0) -- (2.4,2.4)
  node[right, font=\tiny\sffamily] {SNR = 1};

% Data points
\node[marginal] at (0.1536, 0.1059) {};
\node[right, font=\tiny] at (0.18, 0.1059) {modCnt};

\node[noise] at (0.4125, 0.4386) {};
\node[right, font=\tiny] at (0.44, 0.44) {triCnt};

\node[marginal] at (0.2353, 0.2303) {};
\node[right, font=\tiny] at (0.26, 0.23) {edgeCnt};

\node[marginal] at (0.2124, 0.0725) {};
\node[right, font=\tiny] at (0.24, 0.0725) {edgeDens};

\node[marginal] at (0.1465, 0.1249) {};
\node[left, font=\tiny] at (0.13, 0.1249) {meanDeg};

\node[marginal] at (0.4739, 0.2947) {};
\node[right, font=\tiny] at (0.50, 0.29) {maxDeg};

\node[noise] at (0.1244, 0.1373) {};
\node[left, font=\tiny] at (0.10, 0.14) {curvMn};

\node[noise] at (0.2275, 0.2452) {};
\node[left, font=\tiny] at (0.20, 0.26) {curvStd};

\node[marginal] at (0.1633, 0.0660) {};
\node[below, font=\tiny] at (0.1633, 0.04) {curvP50};

\node[marginal] at (0.2717, 0.1949) {};
\node[right, font=\tiny] at (0.30, 0.19) {curvP90};

\node[noise] at (0.5315, 0.5458) {};
\node[right, font=\tiny] at (0.56, 0.55) {core};

\node[marginal] at (0.7827, 0.4689) {};
\node[right, font=\tiny] at (0.81, 0.47) {sat};

\node[marginal] at (2.2302, 1.0875) {};
\node[below, font=\tiny] at (2.23, 1.06) {bridge};

\node[noise] at (0.4767, 0.8034) {};
\node[left, font=\tiny] at (0.45, 0.82) {scaf};

\node[marginal] at (0.8576, 0.7251) {};
\node[right, font=\tiny] at (0.88, 0.73) {bndry};

\end{tikzpicture}
\caption{Signal-to-nuisance landscape.  Points above the SNR=1 line
(gray dashed) carry more noise than signal.  Points above the SNR=3
line (wine dashed) fail the acceptance criterion.  \textbf{All points
are above the SNR=3 line.}  Gold-filled circles are marginal
($1 \leq \text{SNR} < 3$); wine-filled circles have $\text{SNR} < 1$.}
\label{fig:snr-landscape}
\end{figure}

\subsection{The Curvature Shock}
\label{sec:curvature-shock}

Spectral Ecology v9 represented the program's strongest separation
result ($+3.20\%$) and was built entirely on curvature features: the
Ollivier--Ricci curvature distribution was compared using Bhattacharyya
coefficients, with penalties derived from curvature standard deviation
and distortion measures.

\begin{limitation}
\textbf{The features v9 scored have signal-to-nuisance ratios below 1.}
Mean curvature (the primary signal): SNR = 0.91.
Curvature standard deviation (the penalty term): SNR = 0.93.
These are the two features most central to v9's scoring formula.  The
$+3.20\%$ separation was not biological signal---it was structured noise
that happened to correlate with relatedness in the specific coverage
configuration of the CEPH 1463 dataset.
\end{limitation}

Only curvP50 (SNR = 2.47) and curvP90 (SNR = 1.39) achieve ratios above
1, and neither reaches the acceptance threshold of 3.  The curvature
median is relatively stable (self CV = 0.066, the lowest of any feature),
but the cross-sample variation is modest (CV = 0.163), giving a ratio
that is the second-best overall but still insufficient.

This is a devastating finding for the Spectral Ecology program.  The
version that achieved the best separation was, in measurement science
terms, \textbf{scoring noise}.

\subsection{Shuffle Test Results}
\label{sec:shuffle}

Table~\ref{tab:shuffle} presents the results of the read-order shuffle
test.  For each sample, the original module graph (from reads in file
order) is compared to the shuffled graph (same reads, random order).

\begin{table}[htbp]
\centering
\tablestyle
\caption{Shuffle test results.  Read order is randomized with zero change
in biological content.  Percentage changes in graph features are shown.
Any change represents a pure representation artifact.}
\label{tab:shuffle}
\begin{tabular}{@{}l rr rr rr@{}}
\toprule
\rowcolor{table-header}
\textcolor{white}{\textbf{Sample}} &
\multicolumn{2}{c}{\textcolor{white}{\textbf{moduleCount}}} &
\multicolumn{2}{c}{\textcolor{white}{\textbf{triangleCount}}} &
\multicolumn{2}{c}{\textcolor{white}{\textbf{edgeCount}}} \\
\cmidrule(lr){2-3} \cmidrule(lr){4-5} \cmidrule(lr){6-7}
\rowcolor{table-header}
& \textcolor{white}{Orig.} & \textcolor{white}{Shuf.}
& \textcolor{white}{Orig.} & \textcolor{white}{Shuf.}
& \textcolor{white}{Orig.} & \textcolor{white}{Shuf.} \\
\midrule
NA12889 & 457  & 938 \,\textcolor{yalumba-wine}{(+105\%)}
        & 4822 & 4940 \,\textcolor{yalumba-wine}{(+2\%)}
        & 1462 & 2488 \,\textcolor{yalumba-wine}{(+70\%)} \\
NA12890 & 607  & 1023 \,\textcolor{yalumba-wine}{(+69\%)}
        & 1875 & 4650 \,\textcolor{yalumba-wine}{(+148\%)}
        & 1366 & 2377 \,\textcolor{yalumba-wine}{(+74\%)} \\
NA12891 & 712  & 1252 \,\textcolor{yalumba-wine}{(+76\%)}
        & 4867 & 8964 \,\textcolor{yalumba-wine}{(+84\%)}
        & 2251 & 3776 \,\textcolor{yalumba-wine}{(+68\%)} \\
NA12892 & 593  & 1232 \,\textcolor{yalumba-wine}{(+108\%)}
        & 3714 & 10317 \,\textcolor{yalumba-wine}{(+178\%)}
        & 1865 & 3881 \,\textcolor{yalumba-wine}{(+108\%)} \\
NA12877 & 492  & 1395 \,\textcolor{yalumba-wine}{(+184\%)}
        & 1256 & 8875 \,\textcolor{yalumba-wine}{(+607\%)}
        & 1110 & 4045 \,\textcolor{yalumba-wine}{(+264\%)} \\
NA12878 & 671  & 974 \,\textcolor{yalumba-wine}{(+45\%)}
        & 3296 & 5819 \,\textcolor{yalumba-wine}{(+77\%)}
        & 1998 & 2536 \,\textcolor{yalumba-wine}{(+27\%)} \\
\bottomrule
\end{tabular}
\end{table}

\begin{keyfinding}
\textbf{Shuffling read order---a transformation with zero biological
content change---produces catastrophic changes in graph structure.}
Module count increases by 45\% to 184\%.  Triangle count increases by
2\% to 607\%.  Edge count increases by 27\% to 264\%.  The
module-graph representation is \emph{not a function of the genome}.  It
is a function of the genome \textbf{and the read order}.
\end{keyfinding}

\subsection{Per-Sample Shuffle Detail}

The shuffle test reveals striking per-sample variability.  NA12877
(Father) shows the most extreme changes: module count nearly triples
($492 \to 1395$, $+184\%$), triangle count increases sevenfold
($1256 \to 8875$, $+607\%$), and edge count nearly quadruples
($1110 \to 4045$, $+264\%$).

NA12878 (Mother) shows the mildest changes: module count increases by
45\% ($671 \to 974$) and edge count by 27\% ($1998 \to 2536$).  But
even a 27\% change from zero biological perturbation is unacceptable
for a measurement instrument.

\begin{pullquote}[yalumba-wine]
Triangle count for NA12877 changes by \textbf{607\%} when reads are
shuffled.  This is not a small effect.  This is the measurement
equivalent of a thermometer that reads 20$^\circ$C or 140$^\circ$C
depending on which direction you hold it.
\end{pullquote}

The \emph{direction} of the shuffle effect is consistently positive:
shuffled graphs are always larger (more modules, more edges, more
triangles) than the original.  This suggests that file-order reads
produce more locality in k-mer co-occurrence---adjacent reads from the
same genomic region share k-mers that get merged into the same module.
Shuffling breaks this locality, spreading co-occurrences across more
modules and creating more inter-module edges.

\ymarginnote{The asymmetry is telling: shuffling always \emph{increases}
structure.  File order \emph{compresses} the graph by exploiting
genomic locality.}

The asymmetry has a concrete mechanism.  In file-order processing, read
$i$ and read $i+1$ are likely from nearby genomic positions (Illumina
sequencers produce reads in approximately genomic order within a tile).
K-mers shared between adjacent reads are counted as co-occurrences
\emph{within} existing modules, strengthening existing edges.  When
reads are shuffled, the same k-mer pairs co-occur on reads that are
\emph{not} adjacent in processing order, creating new co-occurrence
edges that bridge previously separate graph regions.  The result is more
modules (because threshold-crossing co-occurrences now span more
regions), more edges (because more k-mer pairs cross the co-occurrence
threshold when accumulated from dispersed reads), and dramatically more
triangles (because the bridging edges create short cycles in the
expanded graph).

This mechanism makes the non-determinism even more troubling.  The
representation does not just depend on read order---it depends on
\emph{how reads are spatially distributed on the sequencing instrument}.
Different flow cells, different lanes, and different sequencing runs of
the same sample would produce different module graphs even with
identical biological content.

\subsection{What the Shuffle Test Proves}
\label{sec:shuffle-proves}

The shuffle test establishes a fact that no amount of downstream
sophistication can overcome:

\begin{definition}
A representation is \textbf{deterministic} if it is a function solely
of the biological content of the input.  A representation that changes
when reads are reordered is \emph{non-deterministic}---it depends on
an arbitrary parameter (file order) that carries no biological
information.
\end{definition}

The module-graph representation is non-deterministic.  This has
immediate consequences:

\begin{enumerate}
  \item \textbf{Pairwise comparisons are unreliable.}  If sample $A$'s
    graph depends on the order of $A$'s reads, then the comparison
    $A \leftrightarrow B$ depends on both read orders.  Running the same
    comparison after sorting reads alphabetically would produce a
    different result.
  \item \textbf{No invariant can fix this.}  Eigenvalues,
    curvature distributions, heat kernels, and transport plans are all
    computed from the graph.  If the graph itself is an artifact of read
    order, these invariants inherit the artifact.
  \item \textbf{Self-perturbation CV underestimates true instability.}
    The coverage ladder measures instability under subsampling, which
    preserves read order.  The shuffle test reveals an additional,
    independent source of instability.  The true representation
    variance is the \emph{sum} of both contributions.
\end{enumerate}

\subsection{The Two Least-Bad Features}
\label{sec:least-bad}

Despite the grim overall picture, two features stand out as the
least-bad candidates for a stabilized representation:

\subsubsection{Edge Density (SNR = 2.93)}

Edge density is the ratio of edges to possible edges: $d = 2|E| /
(|V|(|V|-1))$.  It achieves the highest SNR because it
\emph{normalizes} edge count by graph size, partially canceling the
coverage effect.  When coverage drops, both edges and nodes decrease,
and their ratio is more stable than either absolute count.

Self-perturbation CV = 0.0725---the second-lowest of any feature.
Cross-sample CV = 0.2124---moderate but meaningful.

\begin{keyfinding}
Edge density is the only feature within striking distance of the
acceptance threshold.  At SNR = 2.93, it would need only a modest
improvement in self-perturbation stability (e.g., from deterministic
module extraction) to cross 3.0.
\end{keyfinding}

\subsubsection{Curvature Median (curvP50, SNR = 2.47)}

The median Ollivier--Ricci curvature has the lowest self-perturbation
CV of any feature (0.066).  Medians are robust to outliers and
distributional changes, which is why curvP50 is far more stable than
curvMean (0.137) or curvStd (0.245).

The limitation is that curvP50's cross-sample CV is only 0.163---the
variation between genomes is modest.  This is expected: median curvature
captures the ``typical'' edge geometry, which may be more universal
across human genomes than it is genome-specific.

\subsection{Summary of Results}

\begin{keyfinding}
The audit reveals three independent failures of the module-graph
representation:
\begin{enumerate}
  \item \textbf{Coverage instability:} 10 of 15 features are
    MODERATE or UNSTABLE under coverage perturbation.
  \item \textbf{Noise dominance:} 5 of 15 features have
    $\text{SNR} < 1$; no feature reaches $\text{SNR} \geq 3$.
  \item \textbf{Non-determinism:} read-order shuffling changes
    graph structure by 27\% to 607\%.
\end{enumerate}
These failures are \emph{independent}---a representation could be
coverage-stable but order-dependent, or order-independent but
coverage-sensitive.  The current representation fails on all three axes.
\end{keyfinding}


% ════════════════════════════════════════════════════════════════════
\section{Discussion}
\label{sec:discussion}

\subsection{What the Ten-Version Arc Was Actually Optimizing}
\label{sec:actually-optimizing}

The Spectral Ecology program was, in retrospect, optimizing
\emph{invariance to representation noise} rather than \emph{sensitivity
to biological signal}.  Each version's theoretical advance was, in
practice, a new way to compress or average over the unstable parts of
the module graph, hoping that the residual would be biological.

\begin{itemize}
  \item \textbf{v1's eigenvalue density profiles} averaged the spectrum
    of a graph whose topology changed with coverage.  The profiles were
    stable because eigenvalue distributions are naturally smooth, not
    because they captured stable biology.
  \item \textbf{v3's optimal transport} produced spectacular separation
    ($+4.93\%$) because transport distances are sensitive to graph
    structure---including the unstable parts.  The catastrophic gap
    ($-54.67\%$) was the price: transport faithfully detected the
    representation noise.
  \item \textbf{v7's cooperative stability} improved the gap to
    $-0.61\%$ by requiring agreement across multiple invariants.
    Agreement-based scoring is inherently conservative: it suppresses
    outlier features, which happened to be the unstable ones.  This was
    an accidental, implicit feature selection.
  \item \textbf{v9's curvature distributions} achieved the best
    separation ($+3.20\%$) but scored features with SNR below 1.  The
    separation was genuine in the sense that it was reproducible on the
    specific CEPH dataset---but it was not robust to perturbation.
  \item \textbf{v10's hybrid} combined v7 and v9, inheriting v7's
    accidental stability and v9's apparent but fragile signal.
\end{itemize}

The arc's real contribution was \textbf{theoretical}: it developed the
mathematical frameworks (spectral geometry, optimal transport, gauge
theory, Riemannian curvature) that \emph{would} be powerful tools if
applied to a stable representation.  The frameworks themselves are
sound.  The substrate was not.

\begin{limitation}
The reinterpretation is harsh but precise: v7 was the best version not
because it had the best theory, but because its multi-invariant
agreement mechanism \textbf{accidentally} down-weighted the unstable
features.  In measurement science terms, v7 was an ad hoc robust
estimator.  It worked for the wrong reason.
\end{limitation}

This also explains a puzzling pattern in the arc.  More ``theoretically
sophisticated'' versions (v3, v8, v9) often had \emph{worse} gaps than
simpler ones.  The reason is now clear: sophisticated invariants are
more sensitive to graph structure, and sensitivity to structure means
sensitivity to representation artifacts.  Optimal transport (v3) was
the most sensitive method and had the worst gap ($-54.67\%$).  Spectral
summaries (v1) were the least sensitive and had a moderate gap
($-26.87\%$).  The correlation between mathematical power and practical
failure is not a coincidence---it is a direct consequence of the
representation's instability.

\subsection{Why the Spouse Confound Persists}
\label{sec:spouse-confound}

Every version of Spectral Ecology struggled with the spouse pair
(Father--Mother).  The standard explanation was ``similar coverage
produces similar graphs, which inflates pairwise scores.''  This is
correct but incomplete.

\begin{pullquote}[yalumba-wine]
The spouse confound is not a scoring confound.\\It is a
\textbf{representation confound}.
\end{pullquote}

If two samples have similar coverage, their module graphs will have
similar \emph{absolute} structure (similar module counts, similar edge
counts, similar role distributions)---not because their genomes are
similar, but because the graph construction algorithm responds primarily
to the \emph{amount} of data rather than the \emph{content} of data.

No scoring method can fix this.  Null-model deconfounding (v2) subtracts
the expected score under random graphs of the same size, but the size
itself is a coverage artifact.  Self-baseline normalization (v8) divides
by each sample's self-score, but the self-score is also coverage-dependent.
Distributional comparison (v9) compares shapes rather than magnitudes,
but the shapes are coverage-driven.

The fix must be at the representation level:  the graph construction
must be modified so that graphs from the same coverage produce
\emph{systematically different} structures for different genomes.  This
is the canonicalization challenge discussed in
Section~\ref{sec:canonicalization}.

\subsection{The Shuffle Catastrophe}
\label{sec:shuffle-catastrophe}

The shuffle test is the single most important result in this paper.  It
establishes that the module-graph representation is not a function of
the genome---it is a function of the genome \emph{and} an arbitrary
parameter (read order) that carries no biological information.

\ymarginnote{A representation that depends on read order is not
measuring the genome.  It is measuring the sequencer's output buffer.}

The mechanism is clear.  K-mer co-occurrence is computed from reads in
the order they appear.  When reads are in file order (approximately the
order they were sequenced), adjacent reads come from nearby genomic
regions and share k-mers.  This produces dense local co-occurrence that
merges into large, tightly connected modules.  When reads are shuffled,
co-occurrence becomes sparser and more distributed, producing more
modules with different connectivity patterns.

This is a fundamental flaw in the co-occurrence counting step, not a
downstream issue.  The co-occurrence matrix should depend only on
\emph{which} k-mers appear on \emph{which} reads, not on the order in
which reads are processed.  The fact that it does depend on order
reveals that the co-occurrence computation has an implicit windowing or
streaming assumption that makes it path-dependent.

\subsubsection{Magnitude of the Effect}

To put the shuffle effect in perspective:

\begin{itemize}
  \item The \textbf{coverage ladder} reduces data by up to 60\% (from
    50,000 to 20,000 reads).  This causes module count changes of
    roughly 10\% (CV = 0.106).
  \item The \textbf{shuffle test} changes zero data.  It causes module
    count changes of 45\% to 184\%.
\end{itemize}

Read reordering, a biologically vacuous transformation, produces effects
\textbf{4 to 18 times larger} than a 60\% reduction in data.  The
representation is more sensitive to file I/O order than to the amount of
data available.

For triangle count, the effect is even more extreme.  The coverage
ladder produces CV = 0.44.  The shuffle test produces changes up to
607\%.  Triangle count is, for all practical purposes, a random variable
conditioned on read order.

\subsection{Why Role Assignment Is Almost Entirely Noise}
\label{sec:roles-noise}

The five ecological role counts (core, satellite, bridge, scaffold,
boundary) were introduced as a biologically motivated way to
characterize graph structure~\cite{hartwell1999cell}.  The idea was that
roles---nodes classified by their topological function---would capture
higher-order organization that raw graph statistics miss.

The audit shows this idea failed.  All five role counts are UNSTABLE
under self-perturbation (CV = 0.47 to 1.09), and none achieves an SNR
above 2.1.  Scaffold has the worst SNR of any feature (0.59).

The failure mode is straightforward.  Role assignment depends on
thresholds applied to degree, betweenness centrality, and clustering
coefficient.  These three metrics are themselves unstable: degree
varies with coverage, betweenness varies with graph connectivity
(which varies with coverage), and clustering coefficient varies with
triangle count (which varies with both coverage and read order).
Thresholding an unstable metric amplifies instability---nodes near the
threshold flip between roles with small changes in the underlying
metric.

\begin{limitation}
Role assignment, as implemented, is a \textbf{noise amplifier}: it
takes moderate instability in continuous metrics and converts it into
catastrophic instability in discrete categories.  The bridge role is
especially vulnerable because few nodes have the specific combination
of low clustering and high betweenness that defines bridges.  Small
changes in the graph can eliminate all bridges or create many.
\end{limitation}

The lesson is general: \textbf{discretizing continuous features
amplifies instability.}  This is well known in statistics (binning
continuous variables loses information and creates edge effects) but
was not considered when role assignment was
designed~\cite{langfelder2008wgcna}.  A continuous role assignment---where
each node has a probability distribution over roles rather than a hard
assignment---would be more stable, but would also require different
downstream comparison methods.

Consider the concrete example of the bridge role.  In the CEPH data,
bridge counts range from 0 to 459 across samples at 100\% coverage
(Table~\ref{tab:cross-sample}).  The self-perturbation CV is 1.09---the
standard deviation across coverage levels exceeds the mean for several
samples.  For NA12877, the bridge count is 0 at some coverage levels
and nonzero at others, producing a CV that is either undefined (when the
mean is zero) or artificially extreme.  Any algorithm that uses bridge
count as a discriminative feature is, in practice, scoring a random
variable.

Spectral Ecology v2 and v4 weighted ecological role distributions
heavily.  The audit explains their poor performance: they were
effectively adding noise to their scoring function and calling it
``ecological structure.''  The biological metaphor (cores,
satellites, bridges) was appealing but the implementation was
measuring the wrong thing.

\subsection{The Coverage Ladder: What Subsampling Reveals}
\label{sec:coverage-ladder-discussion}

The coverage ladder reveals a scale-sensitivity hierarchy among features:

\begin{enumerate}
  \item \textbf{Scale-invariant} (CV $< 0.10$): edge density, curvP50.
    These are ratios or medians that self-normalize.
  \item \textbf{Scale-sensitive} (CV $0.10$--$0.20$): module count,
    mean degree, curvMean.  These change with coverage but at a
    predictable rate.
  \item \textbf{Scale-dominated} (CV $> 0.20$): everything else.  These
    are dominated by the amount of data rather than its content.
\end{enumerate}

The distinction between scale-invariant and scale-sensitive features
is encouraging.  It suggests that \textbf{normalization is possible}:
if we can identify the scaling law that governs each feature's
dependence on coverage, we can correct for it.

Edge density is naturally normalized (it divides by graph size).  Median
curvature is naturally robust (medians resist outliers).  Both achieve
their relatively good SNRs \emph{because of their mathematical
structure}, not because of anything in the algorithm.  This is a strong
hint: \textbf{features should be designed for scale-invariance from the
start}, not computed as raw counts and then post-hoc normalized.

\subsection{Salvageable Features}
\label{sec:salvageable}

Despite the overall negative result, two features merit further
investigation.

\textbf{Edge density} (SNR = 2.93) is tantalizingly close to the
acceptance threshold.  Its self-perturbation CV of 0.0725 could
plausibly be reduced below 0.05 with deterministic module extraction
(eliminating read-order dependence) and more aggressive frequency
filtering (stabilizing the co-occurrence graph).  If cross-sample CV
remains at 0.21, this would push SNR above 4.

\textbf{Curvature median} (SNR = 2.47) is the most stable feature
by self-perturbation CV (0.066), but its cross-sample CV is only 0.163.
To reach SNR = 3, either self-CV must drop below 0.054 (feasible) or
cross-CV must increase (potentially achievable by using k-mer frequency
distributions rather than binary presence/absence in curvature
computation).

These two features, if stabilized, could form the basis of a minimal
but reliable scoring system.  The key insight is that they are
reliable \emph{because of their mathematical properties} (normalization,
robustness to outliers), not because of any specific biological theory.

\subsection{What Canonicalization Would Need to Achieve}
\label{sec:canonicalization}

To fix the representation, the module extraction pipeline must become
\textbf{canonical}: the output must be a function of the set of reads,
not the sequence of reads.  This requires:

\begin{enumerate}
  \item \textbf{Order-independent co-occurrence.}  Count k-mer
    co-occurrences across all reads simultaneously, not in streaming
    fashion.  This eliminates read-order dependence but requires
    $O(n)$ memory for $n$ unique k-mers.
  \item \textbf{Coverage-matched construction.}  Before constructing
    the graph, subsample each sample to a target coverage level.  This
    eliminates the coverage confound at the representation level.
  \item \textbf{Stable thresholding.}  Replace the current fixed
    threshold (mean $+ 1\sigma$) with a data-adaptive threshold
    that produces stable graph density across coverage levels.
    Rank-based thresholding (keep the top $p\%$ of edges by
    co-occurrence weight) is one candidate.
  \item \textbf{Deterministic module assignment.}  If connected
    components are used, the graph must be connected-component-stable:
    small changes in edge set should not cause large modules to split
    or merge.  Hierarchical clustering with a fixed
    cut~\cite{langfelder2008wgcna} may be more stable than connected
    components.
\end{enumerate}

\begin{definition}
A \textbf{canonical representation} satisfies three properties:
(1)~\emph{determinism}---same read set produces the same graph
regardless of read order; (2)~\emph{scale equivariance}---graph
structure changes predictably with coverage, enabling correction;
(3)~\emph{biological sensitivity}---different genomes produce
different graph structures at the same coverage.
\end{definition}

\subsection{Connection to the Theoretical Framework}
\label{sec:theory-connection}

The theoretical framework paper~\cite{nakahara2003geometry} developed
the concept of ``coarse-graining'' as a functor from raw genomic data
to structural representations:

\begin{equation}
  \Phi: \mathcal{R} \to \mathcal{G}
  \label{eq:functor}
\end{equation}

where $\mathcal{R}$ is the space of read sets and $\mathcal{G}$ is the
space of module graphs.  The framework assumed that $\Phi$ was
\emph{well-behaved}: continuous, approximately surjective within a
genome's equivalence class, and sensitive to biological differences
between genomes.

The audit reveals that $\Phi$ satisfies none of these properties:

\begin{itemize}
  \item It is not continuous: small perturbations of the input (coverage
    changes) produce disproportionate changes in the output (role count
    CVs above 1.0).
  \item It is not well-defined on equivalence classes: reordering
    reads---a transformation within the same equivalence class---produces
    different outputs.
  \item It is not sensitive to biology: for several features
    (curvMean, curvStd, core role, scaffold role), self-perturbation
    variance exceeds cross-sample variance.
\end{itemize}

The theoretical framework's mathematical machinery---spectral geometry,
optimal transport, persistent homology---is not invalidated by this
audit.  But the coarse-graining assumption \emph{is} invalidated.  The
$\Phi$ that the project has been using is not the $\Phi$ that the
framework requires.

Formally, the framework requires that $\Phi$ satisfy three properties
for the downstream comparison functor $\Delta: \mathcal{G} \times
\mathcal{G} \to \mathbb{R}$ to be meaningful:

\begin{enumerate}
  \item \textbf{Well-definedness.}  If $R_1 \sim R_2$ (same read set,
    different order), then $\Phi(R_1) = \Phi(R_2)$.  \emph{Violated:}
    the shuffle test shows $\Phi(R_1) \neq \Phi(R_2)$ by 27--607\%.
  \item \textbf{Lipschitz continuity.}  If $R_1$ and $R_2$ differ by
    a small number of reads, then $d(\Phi(R_1), \Phi(R_2))$ is small.
    \emph{Violated:} removing 60\% of reads changes features by up
    to 173\% (max CV for bridge role).
  \item \textbf{Discriminativeness.}  If $R_1$ and $R_2$ come from
    biologically different genomes, then
    $d(\Phi(R_1), \Phi(R_2))$ is large relative to
    self-perturbation.  \emph{Violated:} SNR $< 3$ for all features.
\end{enumerate}

The theoretical framework assumed these properties held and built
elegant mathematics on top of them.  The audit shows that none holds.
This does not invalidate the mathematics---it invalidates the
\emph{instantiation}.  The functor $\Phi$ must be rebuilt.

\subsection{Implications for Other Structural Genomics Approaches}

This audit is specific to the yalumba module-graph representation, but
its methodology and findings have broader implications.

Any structural genomics method that builds graphs from k-mer
co-occurrence should be subject to the same three-part audit: coverage
ladder, cross-sample variance, and shuffle test.  Methods that skip
the graph step (e.g., Mash~\cite{ondov2016mash}, which computes sketch
distances directly from k-mer sets) may not suffer from the same
non-determinism, but should still be audited for coverage sensitivity.

The signal-to-nuisance ratio framework is general: it can be applied to
any feature set extracted from any representation.  We recommend it as
a standard pre-requisite before publishing pairwise comparison results.

\begin{pullquote}
If you have not measured your representation's self-perturbation
variance, you do not know whether your results are signal or noise.
\end{pullquote}

The shuffle test, in particular, should be adopted universally.  It
costs almost nothing to run (one random permutation per sample) and
catches a class of bugs---order-dependent computation---that is common
in streaming bioinformatics tools.  Any method that gives different
answers when reads are reordered has an implementation-level flaw that
cannot be fixed by downstream processing.

We note that classical methods do not suffer from this flaw.
MinHash~\cite{ondov2016mash} computes sketches from the k-mer
\emph{set}, which is order-independent by construction.
Alignment-based methods~\cite{browning2012ibd, browning2020ibis}
operate on sorted, indexed files where read order has been canonicalized
by genomic coordinate.  The reference-free community should hold itself
to the same standard of order-independence.

\subsection{This Is a Strong Scientific Result, Not a Failure}
\label{sec:not-failure}

Negative results are undervalued in computational
science~\cite{bendavid2006stability}.  This audit is not a failure of
the research program.  It is a \textbf{precise diagnosis} of where the
program's approach breaks down, which algorithms were accidentally
successful and why, and what specific properties a fixed representation
must satisfy.

The audit also \emph{explains} the ten-version arc's trajectory.
The gap narrowed from $-26.87\%$ to $-0.61\%$ not because successive
algorithms got better at detecting biology, but because successive
algorithms got better at \emph{ignoring noise}.  v7's cooperative
stability worked because it implicitly selected the less-noisy
features.  v9 regressed because it explicitly scored the noisier ones.

This understanding is valuable.  It means the program does not need to
abandon its mathematical frameworks.  It needs to rebuild the
representation that those frameworks operate on.

\subsection{Limitations of This Audit}
\label{sec:limitations}

\begin{enumerate}
  \item \textbf{Small sample size.}  Six samples from a single pedigree.
    The cross-sample CV is computed from $n = 6$, which is too small for
    robust variance estimation.  A larger and more diverse sample set
    could produce different SNR values.

  \item \textbf{Single k-mer size.}  All experiments use $k = 21$.
    Different $k$ values may produce more or less stable module graphs.
    The audit should be repeated across a range of $k$ to determine
    whether stability is $k$-dependent.

  \item \textbf{Fixed frequency thresholds.}  The frequency filters
    (min 2, max 100) were chosen heuristically.  Different thresholds
    could change the module graph and potentially improve stability.

  \item \textbf{Single graph construction threshold.}  The edge
    threshold (mean $+ 1\sigma$) is one of many possible choices.
    Rank-based or quantile-based thresholds might produce more stable
    graphs.

  \item \textbf{Linear coverage ladder.}  We tested 100\%, 80\%, 60\%,
    40\%.  Finer coverage steps (e.g., 5\% increments) would reveal
    whether instability is gradual or threshold-driven.

  \item \textbf{Single shuffle.}  We tested one random permutation
    per sample.  Multiple shuffles would yield a distribution of
    shuffle effects, rather than a point estimate.

  \item \textbf{No joint analysis.}  Self-perturbation CV and shuffle
    effects are analyzed independently.  A joint model that
    decomposes total variance into coverage, order, and biological
    components would be more informative.

  \item \textbf{CV as a stability metric.}  The coefficient of
    variation assumes symmetric, approximately normal variation around
    the mean.  For features with zero-inflated distributions (scaffold,
    boundary roles), CV is poorly defined and may overstate instability.
    Intraclass correlation or robust dispersion measures would be
    better suited.

  \item \textbf{No read-length analysis.}  Read length variation could
    interact with co-occurrence counting.  Longer reads produce more
    k-mer pairs per read, potentially changing module structure.  The
    audit does not control for this.

  \item \textbf{SNR threshold is conventional.}  The choice of
    $\text{SNR} \geq 3$ as the acceptance criterion is borrowed from
    psychometrics and is not derived from a formal statistical
    argument specific to genomic comparison.  A different threshold
    could change the conclusions.

  \item \textbf{No comparison to classical methods.}  The audit
    evaluates the module-graph representation in isolation.  It would
    be informative to apply the same perturbation tests to classical
    features (k-mer frequency profiles, MinHash sketches, alignment-based
    statistics) to establish a baseline for ``normal'' feature stability
    in genomics.
\end{enumerate}


% ════════════════════════════════════════════════════════════════════
\section{Recommendations for the Next Phase}
\label{sec:recommendations}

\subsection{Representation Canonicalization}
\label{sec:rec-canonical}

The first priority is eliminating read-order dependence.  The module
extraction pipeline must produce identical output for any permutation
of the same read set.  This requires:

\begin{enumerate}
  \item Compute all k-mer frequencies in a single pass (order-independent
    hash table).
  \item Compute co-occurrence as a function of \emph{which} k-mers
    appear on \emph{which} reads, using read identifiers rather than
    sequential processing.
  \item Threshold co-occurrence using the complete co-occurrence matrix,
    not a streaming approximation.
\end{enumerate}

\ymarginnote{Canonicalization has a memory cost: the full co-occurrence
matrix for $n$ k-mers requires $O(n^2)$ space.  For 50,000 reads at
$k = 21$, this is manageable ($\sim$500K unique k-mers, $\sim$1GB
sparse matrix).}

The expected impact: the shuffle test should produce \textbf{zero change}
in all features.  Any residual change after canonicalization indicates
a bug, not a design issue.

The memory cost is modest.  For 50,000 reads at $k = 21$, we expect
approximately 500,000 unique k-mers.  A sparse co-occurrence matrix
(storing only pairs that co-occur on at least one read) requires
$O(|E|)$ space, where $|E|$ is the number of co-occurring pairs.  In
practice, this is 1--5 million pairs, fitting comfortably in 100--500~MB.
The current streaming approach uses less memory but sacrifices
determinism.  The tradeoff is clear: determinism is non-negotiable;
memory is cheap.

\subsection{Coverage-Matched Subsampling}

Before constructing module graphs for comparison, all samples should be
subsampled to a common coverage level.  The target should be the minimum
coverage across all samples being compared.

\begin{equation}
  n_{\text{target}} = \min_{s \in \mathcal{S}} n_s
  \label{eq:coverage-match}
\end{equation}

where $n_s$ is the read count for sample $s$ and $\mathcal{S}$ is the
set of samples being compared.  This eliminates the systematic
relationship between coverage and graph structure that underlies the
spouse confound.

Multiple random subsamples (e.g., 10 per sample) should be drawn and
averaged, reducing the variance of the subsampled representation.

\subsection{Rank-Based or Quantile-Normalized Features}

Raw feature values are scale-dependent and distribution-dependent.
Replace raw features with their ranks or quantile-normalized values
within the sample:

\begin{equation}
  f_{\text{rank}}(v, s) = \frac{\text{rank}(f(v, s))}
                               {\text{max\_rank}(s)}
  \label{eq:rank-norm}
\end{equation}

This maps all features to $[0, 1]$ regardless of graph size, which
eliminates absolute scale effects.  It also converts unstable absolute
values into more stable relative orderings.

\subsection{Threshold Sweeps}

Instead of using a single edge threshold (mean $+ 1\sigma$), sweep
across thresholds and identify the range where features are most stable:

\begin{equation}
  \tau^* = \arg\min_\tau
    \overline{\text{CV}}_{\text{self}}(f \mid \tau)
  \label{eq:threshold-sweep}
\end{equation}

The optimal threshold may differ by feature.  The sweep should be
performed once (using the self-perturbation test) and the resulting
threshold fixed for all subsequent experiments.

\subsection{The Acceptance Criterion for v11}

\begin{keyfinding}
No new algorithm version should be developed until the following
criterion is met:
\begin{equation}
  \forall\, f \in \mathcal{F}_{\text{scored}}: \quad
  \text{SNR}(f) \geq 3
  \label{eq:acceptance}
\end{equation}
where $\mathcal{F}_{\text{scored}}$ is the set of features used in the
algorithm's scoring formula.  This criterion must be verified on the
CEPH 1463 pedigree using the self-perturbation test (coverage ladder)
and the shuffle test (read-order permutation).  \textbf{Both tests must
yield zero-change results for the shuffle test and SNR $\geq$ 3 for the
coverage ladder.}
\end{keyfinding}

This is a hard gate.  It means that v11 will not be an algorithm
paper---it will be a \textbf{representation paper}.  The algorithm
cannot advance until the substrate is trustworthy.  This is the right
order.

\subsection{Feature Design Principles}

The audit suggests specific principles for designing features that are
likely to achieve acceptable SNR:

\begin{enumerate}
  \item \textbf{Prefer ratios over counts.}  Edge density (SNR = 2.93)
    outperforms edge count (SNR = 1.02) because it normalizes by graph
    size.  Any count-based feature should be divided by the appropriate
    scale factor (number of nodes, number of edges, graph diameter).
  \item \textbf{Prefer medians over means.}  curvP50 (SNR = 2.47)
    outperforms curvMean (SNR = 0.91) and curvStd (SNR = 0.93).
    Medians are robust to the tail behavior that dominates coverage-
    dependent variation.
  \item \textbf{Avoid thresholding.}  Role counts (SNR = 0.59--2.05)
    are uniformly worse than the continuous metrics they are derived
    from.  Keep features continuous whenever possible.
  \item \textbf{Normalize before extraction.}  Rather than extracting
    raw features and normalizing afterward, design the extraction so
    that features are scale-invariant by construction.
  \item \textbf{Test before scoring.}  Every feature proposed for a
    scoring formula must pass the self-perturbation test (CV $< 0.15$)
    and the shuffle test (zero change) \emph{before} it is used in
    pairwise comparison.
\end{enumerate}

\subsection{A Proposed Verification Pipeline}

We propose the following pipeline for validating any new representation
before it is used for algorithmic development:

\begin{enumerate}
  \item \textbf{Shuffle test.}  Permute read order for each sample.
    Require $\Delta f = 0$ for all features.  If any feature changes,
    the representation is non-deterministic and must be fixed before
    proceeding.
  \item \textbf{Coverage ladder.}  Subsample at 100\%, 80\%, 60\%,
    40\%.  Compute self-perturbation CV for each feature.  Classify
    features as STABLE ($< 0.15$), MODERATE ($0.15$--$0.40$), or
    UNSTABLE ($\geq 0.40$).  Only STABLE features proceed.
  \item \textbf{Cross-sample variance.}  Compute cross-sample CV at
    100\% coverage.  Features with cross-sample CV $< 0.05$ have
    insufficient biological variation and should not be scored.
  \item \textbf{Signal-to-nuisance ratio.}  Compute SNR for all
    remaining features.  Only features with $\text{SNR} \geq 3$
    proceed to algorithm development.
  \item \textbf{Algorithm development.}  Only at this stage---after
    the representation is validated---should new scoring methods
    be designed.
\end{enumerate}

This pipeline inverts the order that the project has followed.  For ten
versions, algorithm development preceded representation validation.
Going forward, representation validation gates algorithm development.


% ════════════════════════════════════════════════════════════════════
\section{Conclusion}
\label{sec:conclusion}

This paper reports a negative result that changes the direction of the
yalumba research program.  A systematic perturbation audit of the
module-graph representation---the computational substrate for ten
versions of Spectral Ecology and three prior symbiogenesis
algorithms---reveals three independent failures:

\begin{enumerate}
  \item \textbf{Coverage instability.}  Only 5 of 15 features are
    stable under subsampling.  The remaining 10 features, including all
    ecological role counts, vary more within a sample across coverage
    levels than between different genomes.
  \item \textbf{Noise dominance.}  No feature achieves a
    signal-to-nuisance ratio of 3 or above.  The best ratio is 2.93
    (edge density).  Five features have ratios below 1.0, meaning
    self-perturbation noise exceeds cross-sample variation.
  \item \textbf{Non-determinism.}  Shuffling read order---a
    transformation with zero biological content change---produces
    changes of 27\% to 607\% in graph features.  The representation
    is a function of read order, not just of the genome.
\end{enumerate}

\begin{pullquote}[yalumba-wine]
The ten-version Spectral Ecology arc optimized invariants on an unstable
representation.  The downstream sophistication outran the upstream
identifiability.
\end{pullquote}

These findings explain why the program's gap never turned positive.
The persistent spouse confound (Father--Mother ranking among related
pairs) is not a scoring confound---it is a \textbf{representation
confound}.  Similar coverage produces similar graphs, and no downstream
invariant can distinguish coverage-driven similarity from biological
similarity.

The findings also explain the arc's trajectory.  v7's cooperative
stability worked because it implicitly selected stable features.  v9's
curvature scoring failed (in the SNR sense) because it explicitly
selected unstable ones.  v10's hybrid partially succeeded because it
combined an accidentally stable component (v7) with an accidentally
informative one (v9).

But the program was not wasted.  The theoretical advances are real:
coherence fields, intrinsic curvature distributions, holonomy
deficits, Sinkhorn-regularized transport---these are mathematically
sound frameworks that would be powerful tools for comparing genomic
ecosystems, \textbf{if the ecosystems were reliably
measured}~\cite{chung1997spectral, memoli2011gromov,
edelsbrunner2010computational, peyre2019computational}.

The path forward is clear.  The representation must be
\textbf{canonicalized} (order-independent), \textbf{coverage-matched}
(scale-independent), and \textbf{audited} (SNR $\geq 3$ for every
scored feature) before any new algorithm version is attempted.

This is exactly where real research gets interesting.  The first ten
versions explored the space of invariants.  The next phase explores the
space of representations.  The same mathematical machinery---spectral
geometry, optimal transport, topological persistence---can be applied to
the representation design problem itself: finding graph constructions
whose structure is \emph{invariant} to coverage and read order while
being \emph{sensitive} to biological content.

The audit does not close the book on reference-free genomic ecosystem
comparison.  It opens a harder, more honest, and ultimately more
productive chapter.

The analogy to other fields is encouraging.  In machine learning, the
shift from hand-crafted features to learned representations (the deep
learning revolution) was preceded by a period of negative results showing
that hand-crafted features were unstable and domain-brittle.  In
ecology, the shift from raw species counts to diversity indices was
preceded by demonstrations that raw counts were sampling-effort-
dependent~\cite{bray1957}.  In both cases, the negative results were not
the end of the field---they were the beginning of a more productive
phase.

The same pattern holds here.  The module-graph representation, as
currently implemented, is the ``hand-crafted feature'' of the yalumba
program.  It encodes biological intuition (k-mers that co-occur form
functional modules) but lacks the engineering discipline (determinism,
scale-invariance, measurement validation) needed for reliable
science.  The next phase replaces intuition-driven construction with
measurement-validated construction.

\vspace{1em}

\begin{keyfinding}
\textbf{The most important finding in this paper is not that the
representation is broken.  It is that we now know exactly how it is
broken, exactly which features are salvageable, and exactly what the
acceptance criterion for a fixed representation must be.}  The problem
is now well-posed.  That is progress.
\end{keyfinding}

\vspace{0.5em}

The research thesis of the yalumba project---that human inheritance can
be modeled as the transmission of resilient, interacting genomic
ecosystems---is not invalidated by this audit.  The thesis was never
about module graphs or k-mer co-occurrence.  It was about a level of
biological organization that is richer than haplotype segments and
deeper than token sharing.  The module-graph representation was one
attempt to access that level.  It was not the right attempt, or at
least not a sufficiently careful one.  The audit tells us what ``careful''
means: deterministic, scale-invariant, and validated before use.  The
next representation that satisfies these criteria will be the foundation
for a second decade of algorithmic development---and this time, the
algorithms will be building on rock, not sand.


% ════════════════════════════════════════════════════════════════════
\bibliographystyle{plain}
\bibliography{../references}

\end{document}
